\section{Introduction}\label{sec:intro}

\subsection{A default that is rarely counted}

Two substitutions are treated as safe engineering defaults in deployed learning
systems. The first is structural: when a sequence is long, replace dense
attention by a sparse pattern whose number of attended entries grows more slowly
than $n^2$, and assume that the saving follows. The family is well catalogued
\cite{beltagy2020longformer,zaheer2020bigbird,zhou2021informer,tay2022efficient}:
a sliding window, a handful of global tokens, a block selection, or some
combination of the three. It is standard to call the result ``linear'' or
``sub-quadratic''. The second substitution is representational: when labels are
scarce, replace a task-specific model by a pretrained language model queried
zero-shot or adapted on a few hundred examples, and assume that scale substitutes
for supervision \cite{hegselmann2023tabllm}. Both are attractive for the same
reason. Each replaces an expensive quantity by a cheap one: quadratic work by a
sparsity pattern, labelled data by a prompt.

This paper is about the accounting that both substitutions skip. For attention,
what decides whether a pattern saves anything is the number of attended entries
as a function of $n$, not a qualitative label such as ``sparse'', and once that
number is written down the verdict often reverses. For a language model on a
tabular decision task, what decides whether adaptation helps is the number of
labelled rows; separating that from the interface costs an experiment in which
the interface is repaired while supervision is held fixed. We carry out both
accounts on a single commodity device of the kind these systems are actually
deployed on, rather than on a datacentre accelerator.

\subsection{What we establish}

\paragraph{C1 (entry law, exact)} For the mask that combines a causal sliding
window of width $w$, $s$ attention sinks and a periodic anchor set of period
$g$, the number of attended entries is
\begin{equation}
\lvert \Sset \rvert \;=\; \sum_{i<n}\Big( \min(i{+}1,w) + \min(i{+}1,s)
+ \Big\lfloor \tfrac{i-w}{g} \Big\rfloor_{+} \Big),
\label{eq:count}
\end{equation}
which equals $n(w+s) + (n-w)^2/(2g) + \bigO(n)$ and is verified entry-exact on
the whole grid we measure (Table~\ref{tab:counting}, Figure~\ref{fig:counting}).
Equation~\eqref{eq:count} is not a bound but an identity, so it can be wrong,
and we test it.

\paragraph{C2 (linearity criterion)} Cost is linear in $n$ if and only if the
\emph{global-key budget}, the number of keys visible to every query, is bounded
as $n$ grows (Theorem~\ref{thm:linear}). A fixed period $g$ gives $n/g$ global
keys, so periodic anchoring is quadratic; a fixed set of global tokens gives
$\bigO(1)$ and is linear. The two designs are usually discussed under one
heading. The count separates them, and it is the budget that has to be capped,
not the sparsity.

\paragraph{C3 (granularity law, measured)} Kernels do not execute masks. They
execute blocks. A block-granular kernel materialises a fraction $\min(1, B/g)$ of
the causal key blocks plus a window term, independent of whether the anchor
period is a multiple of the block size, and the decomposition reproduces the
measured density to three decimals at six settings (Table~\ref{tab:block},
Figure~\ref{fig:density}). The consequence is quantitative. Holding the mask
fixed at $g=192$ and coarsening the kernel from $B=32$ to $B=256$ raises the
ratio of executed to advertised density from \BlockInflationGranMin{} to
\BlockInflationGranMax{}$\times$ (Table~\ref{tab:granularity}), and at $B=256$
the kernel materialises every causal block: a pattern advertised as
\BlockElemDensPct\% sparse is executed as dense attention.

\paragraph{C4 (on-device verdict)} On a \GpuMemGiB\,GiB device we measure peak
transient memory and latency for five implementations of the family across
$n\in[256, \LargestN]$. At $n=\LargestN$ the anchor-based variant requires
\VramPswaBig\,MiB against \VramSdpaBig\,MiB for the fused dense kernel
($\VramTimesSdpa\times$) and is \LatencySlowdownVsSdpa$\times$ slower. The same
pattern with the anchors removed is linear in memory and
\LatSlidingRatio$\times$ \emph{faster} than fused dense attention. The cost is
therefore attributable to the anchor set, which is what
Proposition~\ref{prop:count} and Theorem~\ref{thm:linear} predict. We also say
where our implementation is not optimal and how much of the gap is a
kernel-engineering artefact (Section~\ref{sec:attn_disclose}).

\paragraph{C5 (supervision bound)} On identical held-out rows, a $0.5$B
instruction-tuned model reaches AUC-PR \LlmAucPr{} against \GbdtSameAucPr{} for
gradient-boosted trees, a paired gap of \LlmDelta{} (row-level bootstrap
\LlmCiLo--\LlmCiHi; the paired permutation test with \PermSamples{} resamples gives
$p<10^{-3}$). Low-rank adaptation \cite{hu2022lora} on
\SftRows{} rows with per-row continuous targets reduces the JSON parse-failure
rate from \ParseFailZero\% to \SftParseFail\% while AUC-ROC falls from
\LlmAucRoc{} to \SftAucRoc{} and $F_1$ to \SftFOne. Format and discrimination are
therefore separable, and repairing the first does not repair the second. At this
scale what binds is the number of supervision rows. Reading the answer through a
logit interface rather than through generated JSON saves a further factor of
\InterfaceRatio{} in latency, which separates interface cost from model cost.

\paragraph{C6 (protocol)} Each of C1--C5 is a claim that our own apparatus can
refute, and in three cases it did; Section~\ref{sec:audit} describes the
separation between \emph{assertions} (must be true, else the artefact is
inconsistent) and \emph{hypotheses} (may be false, meaning the scientific claim
is refuted), the prohibition on imputing missing measurements, and the
dependency fingerprint that forces a figure to be regenerated when an upstream
result changes. Sections~\ref{sec:attn_withdraw} and~\ref{sec:lang} report the
refuted claims in full.

\subsection{Why this is a contribution to artificial intelligence}

The relevance does not come from tuning one fraud pipeline. It comes from the
fact that three quantities predict whether two of the most widely deployed
substitutions in applied AI will help or hurt: the global-key budget of an
attention pattern, the block inflation factor of the kernel that executes it, and
the number of labelled rows per decision class. Each is an integer or a ratio
available before training starts, and each is falsifiable by an experiment that
fits on one GPU. We make the predictions, we test them, and
in two cases we find the opposite of what the current literature's vocabulary
suggests. Since the same vocabulary (``sparse is cheaper'', ``pretrained models
generalise'') is what practitioners use to choose an architecture under a memory
or label budget, a correction to the vocabulary is worth more than another point
of AUC-PR on one benchmark.

The work touches three parts of AI research, though not evenly. Attention has
recently been analysed as an object with a capacity and a structure of its own
\cite{baldi2023quarks}; what we add is the complementary accounting of what a
mask costs to execute, as opposed to what it can represent. The empirical
literature on when learned models beat strong tabular baselines
\cite{gorishniy2021revisiting,grinsztajn2022why} supplies the second comparison,
which we extend to pretrained language models under an explicit cost and
supervision budget and split into the two failure modes that are usually
conflated. Cost-sensitive decision making \cite{elkan2001foundations} supplies
the criterion, kept unchanged throughout, which is what makes the efficiency
question a decision question rather than a benchmark question, and why our
thresholds are set on validation data and evaluated once.

We also take a position on method. A negative result about a widely used default
is only useful if the community can check it, so the paper is written against an
artefact in which every number is regenerated from a stored result file at build
time, every claim carries a status, and blocked measurements are listed with the
command that would unblock them rather than replaced by an estimate.
Section~\ref{sec:protocol} describes the apparatus and \ref{app:artifact} gives
the reproduction commands. That the apparatus refuted three of our own
intermediate claims is, to us, evidence that it works.

\subsection{Scope and what we do not claim}

We do not claim that sparse attention is useless, nor that language models
cannot make tabular decisions. We claim that one specific and popular mask
family, with a fixed anchor period executed at block granularity on a
memory-constrained device, is dominated by a fused dense kernel at every length
that device can hold; and that a $0.5$B model adapted on $10^2$ rows is
supervision-limited rather than interface-limited. Where the evidence stops we
say so: we measure two model sizes, not a scale law; we measure one GPU, one
EVM and two datasets; the attestation and proof-carrying-inference paths that
motivated the original system are not measured at all and are reported as
blocked (Section~\ref{sec:system}), not as estimated.

\subsection{Organisation}

Section~\ref{sec:setting} fixes notation and the decision problem.
Section~\ref{sec:related} relates the work to efficient attention, cost-sensitive
learning and tabular language modelling, and states the closest prior work and
the difference. Section~\ref{sec:theory} proves the two laws and derives the
hypotheses tested later. Section~\ref{sec:protocol} describes the measurement
protocol. Sections~\ref{sec:attention}, \ref{sec:lang} and \ref{sec:system}
report the attention, language-model and end-to-end results.
Section~\ref{sec:discussion} discusses implications, threats to validity and the
limits of the evidence. \ref{app:proofs} contains the proofs and
\ref{app:artifact} the reproduction recipe.
